\subsubsection{Park Formalization}\label{ssub:park}
\textit{(Joao, C)}\quad
The advanced allocation method is inspired by Park and Song's idea of using
predictive process monitoring for prescriptive resource allocation
decisions~\cite{park2023predictiveallocation}. Their work combines predictions
about running cases with an optimization step for non-clairvoyant process
instances. The implementation in this simulator keeps the central idea, namely
using predicted future work to influence current resource allocation, but it uses
a lightweight assignment heuristic instead of reproducing the full minimum-cost
maximum-flow formulation.

At each allocation epoch, the method constructs two candidate sets. The first set
contains currently enabled waiting tasks, which can be assigned immediately to
available authorized resources. The second set contains predicted future tasks.
Predicted candidates cannot be executed immediately; they can only create a
reservation, meaning that a compatible resource is intentionally kept idle for a
predicted case/activity pair. Before a prediction is consumed by the allocator,
the integration layer filters it against the currently BPMN-valid candidate
activities. Thus, reservations are restricted to process-model-compatible future
work.

For a resource $r$ and candidate $t$, the pairwise assignment cost is
\[
C(r,t) = \hat{p}(r,t) - \alpha_w\,w(t) - \alpha_\pi\,\pi(t) + \rho(r,t),
\]
where $\hat{p}(r,t)$ is the estimated processing time, $w(t)$ is the current
waiting time and $\pi(t)$ is the task priority. The first term favors
resource-task pairs with shorter expected processing time, while the negative
waiting and priority terms make already-waiting or important tasks more
attractive. For currently enabled tasks, the prediction penalty $\rho$ is $0$.
For predicted future candidates,
\[
\rho(r,t) = \lambda_u(1-P(t)) + \lambda_n(1-P(t)) + \lambda_d d(t),
\]
where $P(t)$ is the prediction probability, $d(t)$ is the expected delay,
$\lambda_u$ controls uncertainty aversion, $\lambda_n$ is the no-show penalty and
$\lambda_d$ controls the penalty for reserving a resource for a distant future
task. The canonical run normalizes time costs by one hour, uses a no-show weight
of $1.0$ and disables the additional future-delay term. This makes strategic
idling possible but conservative: uncertain predictions must compensate their
penalty before they can dominate already waiting work.

Rather than selecting independently for each resource, the implementation solves
a global assignment problem for the current allocation epoch. Rows correspond to
available resources and columns correspond to current or predicted candidates.
Infeasible resource-candidate pairs are masked out, each candidate can be used at
most once and dummy idle columns are available as fallback options. The allocator
also exposes a one-step temporal lookahead extension: a currently enabled task
can receive an additional penalty if assigning it would keep the resource busy
beyond the expected arrival time of a compatible predicted task inside the
planning horizon.

This method should therefore be interpreted as a simulator-compatible,
rolling-epoch approximation inspired by Park and Song, not as a complete
reimplementation of the original optimization model. It does not solve a
multi-period MILP or stochastic min-cost flow and it does not construct a full
future schedule. Its quality depends on the prediction stream, calibration,
train-split processing-time estimates, reservation thresholds and release rules.
These empirical effects are discussed only in the evaluation section.
